\documentclass[10pt,a4paper]{amsart}
\usepackage[margin=19mm]{geometry}
\usepackage{amsmath,amssymb,mathtools}
\usepackage[scr=rsfs,cal=euler]{mathalfa}
\usepackage{xfrac}
\usepackage[unicode,hidelinks,bookmarks=false]{hyperref}
\newcommand{\ie}{i.e.\ }
\newcommand{\cf}{cf.\ }
\newcommand{\resp}{respectively\ }
\newcommand{\Subsection}{\subsection}
\providecommand{\nobreakdash}{\nobreak\hbox{-}}
\setlength{\emergencystretch}{2em}
\multlinegap0pt
\usepackage{bm}
\usepackage{bbm}
\usepackage{dsfont}
\usepackage{xspace}
\usepackage{cancel}
\usepackage{mathtools}
\usepackage{paralist}
\usepackage{amsthm}
\makeatletter
\@ifundefined{theorem}{\newtheorem{theorem}{Theorem}}{}
\@ifundefined{lemma}{\newtheorem{lemma}[theorem]{Lemma}}{}
\makeatother
\newcommand{\RR}{\mathbb{R}}
\newcommand{\ZZ}{\mathbb{Z}}
\newcommand{\abs}[1]{\left|#1\right|}
\newcommand{\ceil}[1]{\left\lceil #1 \right\rceil}
\newcommand{\e}{\varepsilon}
\newcommand{\floor}[1]{\left\lfloor #1 \right\rfloor}
\newcommand{\fpa}[1]{\left\lVert #1 \right\rVert_{\RR/\ZZ}}
\renewcommand{\subset}{\subseteq}
\title[Decidability of Extensions of Presburger Arithmetic by Hardy]{Decidability of Extensions of Presburger Arithmetic by Hardy Field Functions}
\author{Hera Brown, Jakub Konieczny}
\date{Source: arXiv:2508.19206v1 (CC BY 4.0)}
\begin{document}
\maketitle

\noindent\emph{Source and licence.} Decidability of Extensions of Presburger Arithmetic by Hardy Field Functions. Version arXiv:2508.19206v1 (CC BY 4.0), distributed under the Creative Commons Attribution 4.0 International licence, as recorded in the arXiv licence metadata for this exact version and shown on the abstract page https://arxiv.org/abs/2508.19206v1. Full rights notice: licenses/arxiv-2508.19206-CC-BY-4.0.txt.\par\smallskip

\noindent\textit{Editorial scope.} Counted unit: the Lemma whose source label is \texttt{lem:lambda-sufficient} in the article (source0.tex lines 635--645, source label \texttt{lem:lambda-sufficient}), delivered together with the article's own complete original proof of it (source0.tex lines 646--691, one \texttt{proof} environment of 46 lines). No mathematics is written, repaired or re-derived: statement and proof are the article's own text line for line. Required context retained uncounted: none. Every definition, display and label of those lines is retained unchanged. Omitted from this package and not redistributed: 1--634, 692--793. Nothing inside a retained excerpt is omitted or altered: every retained body is delivered exactly as the article's lines are written, so no omission receipt of any kind accompanies this package. Citations: the retained excerpts carry no citation command at all, so no bibliography entry of the article is transcribed and no reference is redistributed. Reference adaptation: none was needed and none was made; every reference inside the delivered unit resolves inside it, so no unresolved reference or citation is produced. Printed slips kept literal: no printed word, symbol, hypothesis or number of the counted statement or of its proof was altered. Layout and class: the article's own class is \texttt{amsart} with packages including amsmath, amsthm, listings, xcolor, mathtools, paralist, cite, thmtools, thm-restate, latexsym, textcomp, amssymb, enumitem, hyperref; the wrapper is the lane's pinned \texttt{amsart} editorial wrapper at 10pt on A4, which restates the theorem-like environments the retained text needs and adds the article's own macro definitions that the retained text needs (\texttt{\textbackslash RR}, \texttt{\textbackslash ZZ}, \texttt{\textbackslash abs}, \texttt{\textbackslash ceil}, \texttt{\textbackslash e}, \texttt{\textbackslash floor}, \texttt{\textbackslash fpa}, \texttt{\textbackslash subset}), with extra wrapper packages (bm, bbm, dsfont, xspace, cancel, mathtools, paralist). The delivered numbering is the unit's own. Assets: no retained span contains a figure environment, a graphics-inclusion command, a PDF-page inclusion, a table or a TikZ picture; no image file is copied into this package, the package declares no asset dependency and no image file is redistributed with it. Licence: version arXiv:2508.19206v1 of this article is distributed under the Creative Commons Attribution 4.0 International licence, as recorded in the arXiv licence metadata for this exact version and shown on the abstract page https://arxiv.org/abs/2508.19206v1. A full rights notice is delivered at licenses/arxiv-2508.19206-CC-BY-4.0.txt. Characters: every retained excerpt is pure ASCII, so the delivered engine, XeLaTeX with its default Latin Modern text font, reports no missing character.
\begin{lemma}
\label{lem:lambda-sufficient}

	Let $f \colon [x_0,\infty) \to \RR$ be a Hardy field function with $x \prec f(x) \prec x^2$. Then for any real $x_1 \geq x_0$ and $\e_0,\e_1,\e_2 > 0$ there exists integer $n_0$ such that for each $n \geq n_0$ there exists integer $N$ satisfying the following conditions:
\begin{inparaenum}
\item $N \geq x_1$;
\item $\fpa{f(N)} < \e_0$;
\item $\abs{f'(N) - n} < \e_1$;
\item $f''(N) < \e_2$.
\end{inparaenum}
\end{lemma}

\begin{proof}
	Decreasing $\e_0,\e_1,\e_2$ if necessary (in this order), we may freely assume that $\e_0,\e_1,\e_2 < 1/10$ and we have the inequalities:
\begin{align*}
	\e_1 &< \frac{10}{3} \e_0,&
	\frac{\e_1}{\e_2} &> \frac{25}{\e_1} + 5 > 50.
\end{align*}
(The motivation behind these inequalities will become apparent in the course of the argument). Picking a sufficiently large value of $n_0$ we may freely assume that there exists $X_0 \geq x_1$ such that $f'(X_0) = n$, $f''(X_0) < \e_2$, and $f,f',f''$ are strictly monotone on $[X_0,\infty)$. Let $X_1,X_2,X_3$ be specified by
	\begin{align*}
	f'(X_1) &= n + \frac{1}{5} \e_1,&
	f'(X_2) &= n + \frac{2}{5} \e_1,&
	f'(X_3) &= n + \frac{3}{5} \e_1.&	
	\end{align*}
	Let also $N_1 = \ceil{X_1} < X_1 + 1$ and $N_3 = \floor{X_3} > X_3 - 1$. 
	Note that $f'(X_1) - f'(X_0) = \e_1/5$ and $f''(x) \leq f''(X_0) \leq \e_2$ for $x \in [X_0,X_1]$, so the mean value theorem implies that
	\begin{align*}
	X_1 - X_0 \geq  \frac{\e_1}{5\e_2} \geq 10.
	\end{align*}		
	By the same token, we have $X_2-X_1, X_3-X_2 \geq 10$, which in particular implies that $N_3-N_1 \geq 10$ (which is not strictly speaking necessary, but ensures that all the points under consideration appear in the order we expect). 	
	For all integers $N$ with $N_1 \leq N \leq N_3$ we have:
\begin{align*}
N &\geq x_1,&
n+\frac{1}{5}\e_1 &\leq f'(N) \leq n+\frac{3}{5}\e_1,&
0 < f''(N) \leq \e_2,&
\end{align*}	
which implies that $N$ satisfies three out of the four desired conditions.

It remains to show that there exists $N$ with $N_1 \leq N \leq N_3$ such that $\fpa{f(N)} < \e_0$. We will in fact show that for each arc $A \subset \RR/\ZZ$ with length $2\e_0$ there exists $N$ with $N_1 \leq N \leq N_3$ such that $f(N) \bmod 1 \in A$. Consider the function $\tilde f(N) = f(N) - nN$. Since $\tilde f(N) \equiv f(N) \bmod 1$, in the previous condition we may equally well require that $\tilde f(N) \bmod 1 \in A$. By the mean value theorem, for $N_1 \leq N < N_3$, the value of the expression 
\begin{align*}
	\tilde f(N+1) - \tilde f(N) = f(N+1) - f(N) - n 
\end{align*}
belongs to the interval $\left[\frac{1}{5} \e_1, \frac{3}{5} \e_1 \right]$.
Since $3\e_1/5 < 2\e_0$, it follows that, for each interval $I$ of length $2\e_0$ contained in the interval $[f(N_1),f(N_3)]$, there exists some $N$ with $N_1 \leq N \leq N_3$ with $\tilde f(N)$ intersecting $I$. By the mean value theorem, we have 
\begin{align*}
	N_3 - N_1 \geq X_3 - X_1 - 2 &\geq 
	\frac{f'(X_3) - f'(X_1)}{\max_{X_1 \leq x \leq X_3} f''(x) } - 2 \\
	&= \frac{f'(X_3) - f'(X_1)}{ f''(X_1) } - 2 \geq  \frac{2\e_1}{5\e_2} - 2 \geq \frac{10}{\e_1}.
\end{align*}
 
As a consequence, using the mean value theorem yet again, we get
\begin{align*}
	f(N_3) - f(N_1) &\geq (N_3-N_1) \min_{N_1 \leq x \leq N_3} f'(x) 
	\\ & = (N_3-N_1) f'(N_1) 
	\geq \frac{10}{\e_1} \cdot \frac{\e_1}{5} = 2.
\end{align*}
	Thus, we can find an interval $\tilde A \subset [f(N_1),f(N_3)]$ such that $\tilde A \bmod 1 = A$, and we can find $N$ with $N_1 \leq N \leq N_3$ with $\tilde f(N) \in \tilde A$. This completes the argument.
\end{proof}

\end{document}
